\documentclass[11pt,a4paper]{article}
\usepackage[T1]{fontenc}
\usepackage[utf8]{inputenc}
\usepackage[ngerman]{babel}
\usepackage[margin=2.5cm]{geometry}
\usepackage[hidelinks]{hyperref}
\usepackage{array}
\setlength{\parindent}{0pt}
\setlength{\emergencystretch}{2em}
\renewcommand{\arraystretch}{1.25}
\begin{document}
\shorthandoff{"}
\begin{center}
  {\Large\textbf{Webtechnologien -- Übungsblatt 4}}
\end{center}
\section*{Aufgabe 4: JavaScript Prototypes}
\textbf{Aufgabe:} Ein Bild-Objekt in JSON-Schreibweise mit
\texttt{computeSize()}, danach ein Konstruktor \texttt{Image} mit
\texttt{print()}, eine Methode \texttt{printMore()} über den Prototype und
ein Konstruktor \texttt{Video}, der von \texttt{Image} erbt.

\medskip
\textbf{Ergebnis:} \texttt{image.js}, eingebunden in \texttt{index.html};
die Ausgabe steht in der Konsole (F12). Für Aufgabe 4 gibt es keine
Vorlage in der Angabe, die ganze Datei ist also neu; jede Teilaufgabe steht
zwischen \texttt{[NEU] Aufgabe 4x} und \texttt{[ENDE NEU]}.
\begin{itemize}
  \item \textbf{a)} \texttt{img} ist ein Objektliteral mit
    \texttt{width: 1920}, \texttt{height: 1080}, \texttt{bitdepth: 24} und
    der Methode \texttt{computeSize()}. Sie setzt
    \texttt{this.pixels = width * height} und
    \texttt{this.rawsize = pixels * bitdepth / 8} (Bits durch 8 = Bytes).
    Ausgabe: 2073600 Pixel, 6220800 Bytes.
  \item \textbf{b)} \texttt{function Image(width, height, bitdepth)} setzt
    die drei Werte und rechnet \texttt{pixels} und \texttt{rawsize} sofort
    aus. \texttt{print()} wird im Konstruktor mit
    \texttt{this.print = function () \{...\}} angelegt und gibt mit
    \texttt{for (let key in this)} alle Eigenschaften aus; Funktionen werden
    mit \texttt{typeof} übersprungen, damit nur Daten erscheinen.
    Testobjekte: 800x600, 1080x1920, 500x500.
  \item \textbf{c)} \texttt{Image.prototype.printMore} ruft
    \texttt{this.print()} auf und gibt danach \glqq landscape\grqq{}
    (breiter als hoch), \glqq portrait\grqq{} (höher als breit) oder
    \glqq square\grqq{} aus. Die Objekte aus b) können die Methode nutzen,
    obwohl sie vorher gebaut wurden.
  \item \textbf{d)} Unterschied: \texttt{print} steht im Konstruktor, also
    hat \emph{jedes} Objekt eine eigene Kopie der Funktion.
    \texttt{printMore} steht im Prototype, ist nur \emph{einmal} im
    Speicher und wird von allen Objekten geteilt. Beide lassen sich zur
    Laufzeit durch Zuweisung ändern: \texttt{small.print = function ...}
    ändert nur \texttt{small}, \texttt{tall.print} bleibt gleich.
    \texttt{Image.prototype.printMore = function ...} ändert die Methode
    sofort für alle Bilder und Videos, auch für schon erzeugte.
  \item \textbf{e)} \texttt{function Video(width, height, bitdepth,
    duration, framerate)} ruft zuerst \texttt{Image.call(this, width,
    height, bitdepth)} auf, damit Breite, Pixel, \texttt{print} usw. auch am
    Video gesetzt werden, und speichert dann \texttt{duration} und
    \texttt{framerate}. Mit
    \begin{center}\texttt{Video.prototype = Object.create(Image.prototype)}\end{center}
    erbt der Prototype von
    \texttt{Image.prototype}; \texttt{constructor} wird wieder auf
    \texttt{Video} gesetzt. \texttt{totalFrames()} liegt im Prototype von
    \texttt{Video} und gibt \texttt{duration * framerate} zurück.
  \item \textbf{f)} Ja. \texttt{clip.printMore()} wird zuerst am Objekt,
    dann in \texttt{Video.prototype} und dann in \texttt{Image.prototype}
    gesucht (Prototype-Kette) und dort gefunden. Ohne
    \texttt{Object.create} würde das nicht gehen.
\end{itemize}

\textbf{Hinweis:} Teil d) wird ganz am Ende des Scripts vorgeführt, damit
die Ausgaben von b), c), e) und f) noch die ursprünglichen Funktionen
verwenden. Der eigene Konstruktor \texttt{Image} überdeckt im Browser den
eingebauten \texttt{Image}-Konstruktor für \texttt{<img>}-Elemente; das
stört hier nicht, weil die Seite ihn nicht verwendet.

\medskip
\textbf{Test:} In Chrome, \texttt{index.html} geöffnet, Konsole:
\begin{center}
\begin{tabular}{|>{\raggedright\arraybackslash}p{4.6cm}|>{\raggedright\arraybackslash}p{9.1cm}|}
\hline
\textbf{Aufruf} & \textbf{Ausgabe} \\ \hline
\texttt{img.computeSize()} & pixels 2073600, rawsize 6220800 bytes \\ \hline
\texttt{new Image(800, 600, 8).print()} & width 800, height 600,
  bitdepth 8, pixels 480000, rawsize 480000 \\ \hline
\texttt{printMore()} für 800x600, 1080x1920, 500x500 & zusätzlich
  landscape, portrait, square \\ \hline
\texttt{new Video(1280, 720, 24, 10, 25)} & \texttt{print()} zeigt auch
  duration 10 und framerate 25, \texttt{totalFrames()} = 250 \\ \hline
\texttt{new Video(1080, 1920, 24, 60, 30)} & \texttt{totalFrames()} = 1800,
  \texttt{printMore()} $\rightarrow$ portrait \\ \hline
d) neue Zuweisungen & \glqq small: 800x600\grqq, \texttt{tall.print()}
  unverändert, \glqq printMore replaced: \dots\grqq{} bei Bild und Video \\ \hline
\end{tabular}
\end{center}
Keine Fehler in der Konsole; mit \texttt{node image.js} gleiche Ausgabe.

\medskip
{\small \textbf{Hilfsmittel:} Claude Code (Claude Opus 5.5) für Erstellung, Prüfung und Dokumentation.}
\end{document}
